% !TEX root = ../main.tex
\section{Cómo se calcula la evolución}
\label{sec:metodo-plano}

\subsection{Aproximar el campo y sus derivadas}
El programa almacena el campo en $N$ posiciones igualmente separadas.
La distancia entre dos posiciones consecutivas es $h=\widehat L/N$.
Se almacena el extremo izquierdo y se omite la copia del extremo derecho,
porque ambos representan el mismo punto del dominio periódico. Para
calcular junto a un extremo se utiliza como vecino el punto correspondiente
del extremo opuesto.

La ecuación necesita una segunda derivada espacial, pero el programa
dispone de valores en posiciones separadas. La aproximamos utilizando el
valor de un punto y los de sus dos vecinos. Su origen se ve al desarrollar
por Taylor una función suave $f$:
\[
\begin{aligned}
f(x+h)&=f(x)+hf'(x)+\tfrac12h^2f''(x)+\cdots,\\
f(x-h)&=f(x)-hf'(x)+\tfrac12h^2f''(x)+\cdots.
\end{aligned}
\]
Al sumar, los términos de la primera derivada se cancelan. Despejar la
segunda derivada conduce a
\[
f''(x)\simeq\frac{f(x+h)-2f(x)+f(x-h)}{h^2}.
\]
Esta fórmula compara el valor central con los de alrededor y aproxima
cómo cambia la pendiente del perfil. Se llama diferencia centrada porque
utiliza un vecino a cada lado. Para una función suficientemente suave,
el primer error que queda después de la cancelación es proporcional a
$h^2$ al hacer $h$ pequeño. En los perfiles sinusoidales utilizados se
cumple esa regularidad.

\subsection{Avanzar de un instante al siguiente}
El tiempo también se divide en pasos de tamaño $\Delta\widehat t$.
Escribimos $\Phi_i^n$ para el valor numérico del campo adimensional en
la posición $i$ y en el instante $n\Delta\widehat t$. El subíndice $i$
identifica una posición; el superíndice $n$ identifica un instante, no
una potencia.

Aproximamos la segunda derivada temporal del mismo modo, utilizando el
instante anterior, el actual y el siguiente. Al sustituir ambas
aproximaciones en Klein--Gordon y despejar el valor siguiente resulta
\begin{equation}
\begin{aligned}
\Phi_i^{n+1}={}&2\Phi_i^n-\Phi_i^{n-1}\\
&+(\Delta\widehat t)^2
\left(\frac{\Phi_{i+1}^n-2\Phi_i^n+\Phi_{i-1}^n}{h^2}
-\widehat\mu^2\Phi_i^n\right).
\end{aligned}
\label{eq:esquema-plano}
\end{equation}
La expresión entre paréntesis aproxima la segunda derivada temporal que
predice la EDP. Así, cada paso utiliza el campo actual, sus vecinos y el
campo del instante anterior. La regla es explícita: todos esos valores
ya se conocen cuando se calcula el siguiente.

Para empezar necesitamos dos instantes. El primero se fija con el perfil
inicial. El segundo se aproxima mediante Taylor en el tiempo. Como la
velocidad inicial es cero, el término lineal desaparece y queda
\begin{equation}
\Phi_i^1=\Phi_i^0+\frac{(\Delta\widehat t)^2}{2}
\left(\frac{\Phi_{i+1}^0-2\Phi_i^0+\Phi_{i-1}^0}{h^2}
-\widehat\mu^2\Phi_i^0\right).
\label{eq:inicio-plano}
\end{equation}
Esto incorpora el cambio inicial que exige la ecuación aunque el campo
parta con velocidad nula. A partir de esos dos instantes se repite
\eqref{eq:esquema-plano} hasta el tiempo final.

\subsection{Elección del paso y comprobación del error}
El paso temporal debe elegirse junto con la separación espacial. Un paso
demasiado grande puede hacer crecer errores del cálculo aunque la
solución exacta permanezca acotada. Para este esquema periódico utilizamos
la restricción de estabilidad
\begin{equation}
\Delta\widehat t<\frac{2}{\sqrt{\widehat\mu^2+4/h^2}}.
\label{eq:estabilidad-plana}
\end{equation}
Es una condición del procedimiento numérico. En las ejecuciones se propone
un paso igual a $0.8$ veces esa cota y se reduce ligeramente, si hace falta,
para que un número entero de pasos termine exactamente en
$\widehat t=20$. Al aumentar $N$ también disminuye el paso temporal.

Las diferencias centradas aproximan ambas derivadas con errores de orden
dos para soluciones suaves. Por ello esperamos que el error disminuya
aproximadamente por un factor cuatro al reducir a la mitad los pasos
espacial y temporal. Esta expectativa se contrasta con las cuatro mallas
utilizadas; los resultados y el alcance de esa comprobación se explican
a continuación.
